\chapter{Thermodynamic twin of the CFM56-7B26}
\label{ch:model}

This chapter documents the completed part of phase F1: the CFM56-7B26 performance model in
pyCycle~\cite{Hendricks2019}, its design point (work package WP1.1) and its off-design
calibration against certified measured data (WP1.2). All result figures and tables are produced
by the evidence-generation script (\cref{sec:reproducibility}).

\section{What a cycle model is}

A 0-D \emph{cycle model} represents the engine as a chain of thermodynamic components ---inlet,
fan, compressors, combustor, turbines, nozzles--- each transforming the air's pressure,
temperature and composition according to its governing equations, without resolving any
internal geometry. Each turbomachine carries a \emph{component map}: a table, obtained from rig
tests or scaled from a similar machine, giving its efficiency and flow as a function of
rotational speed and operating point. The model runs in two modes:

\begin{description}
  \item[Design mode] sizes the engine: given targets (thrust, component pressure ratios,
    turbine inlet temperature), it computes the flow areas, map scale factors and shaft-work
    balance that a machine meeting those targets must have.
  \item[Off-design mode] answers the service question: given the now-fixed geometry, what does
    the engine do at any other flight condition and thrust setting? This is the mode the EHM
    baseline decks and the influence coefficient matrix are built from.
\end{description}

pyCycle solves both modes with a Newton solver: it iterates on a set of unknowns (mass flow,
fuel--air ratio, spool speeds, map operating points) until a set of residual equations
(nozzle-area match, shaft power balance, thrust target) is driven to zero. Newton solvers are
powerful but fragile far from a good initial guess ---a fragility that shaped several decisions
below.

\section{Strategy: adapt, don't build from scratch}
\label{sec:morph}

pyCycle ships a validated example of a two-spool, separate-flow, cooled high-bypass turbofan
(\code{high\_bypass\_turbofan.py}) with the full component graph, the turbine-cooling bleed
network, shaft balances and off-design solver wiring already working. Building the CFM56 from a
blank file would re-derive all of that plumbing and typically dies in Newton divergence.
Instead, a \emph{morphing} protocol is adopted: the example is vendored into the repository as
a known-good regression reference, and the model is steered toward the CFM56-7B26 by changing
\textbf{at most three parameters per step}, requiring convergence at every step before the
next. Every converged step is a commit; if a step diverges, the change is bisected.

All design decisions were fixed in writing \emph{before} coding (\code{docs/f1-model-spec.md}):
the -7B26 variant; tabular thermodynamics for speed with a CEA cross-check planned (CEA is
NASA's chemical-equilibrium code ---more accurate, roughly an order of magnitude slower); the
design point at cruise, Mach~0.78 / \SI{35000}{ft} / ISA (the ``aerodynamic design point''
convention: an engine is designed for cruise, where it spends most of its life, and takeoff is
an off-design condition); pyCycle's generic component maps, scaled ---no public CFM56 maps
exist, a declared limitation; thrust set through corrected N1 in off-design (the -7B is an
N1-rated engine); pyCycle-native units inside the model with SI only at the boundary; and a
table of anticipated solver failure modes, each with a planned countermeasure ---one of which
proved prophetic (\cref{sec:idle}).

\section{Design point (WP1.1)}
\label{sec:design}

The cycle graph comprises: flight conditions, inlet, fan, flow splitter, booster, HPC with
three bleed ports (HPT vane and blade cooling, customer bleed), a post-compressor bleed feeding
the LPT, combustor, cooled HPT and LPT, ducts with pressure losses, convergent core and bypass
nozzles, the two shafts, and an accessory power extraction of \SI{250}{hp} on the HP shaft
(driving the aircraft's generators and pumps). In design mode, four implicit balances close the
cycle: the inlet mass flow $W$ meets the thrust target; the fuel--air ratio (FAR) meets the
target turbine inlet temperature $T_{t4}$; and the HPT and LPT expansion ratios zero the net
power on each shaft (each turbine produces exactly what its compressors and losses consume).

\begin{table}[htbp]
  \centering\small
  \caption{Design-point results (Mach 0.78, \SI{35000}{ft}, ISA). Auto-generated from the
  model.}
  \label{tab:design}
  % AUTOGENERATED by scripts/make_report_assets.py — do NOT edit by hand
\begin{tabular}{ll}
\toprule
Quantity & Value at the design point \\
\midrule
  Net thrust $F_n$ & \SI{24.38}{\kilo\newton} (\num{5480} lbf) \\
  TSFC & \num{0.6214} lb/(lbf$\cdot$h) \\
  Bypass ratio (BPR) & \num{5.10} \\
  Overall pressure ratio (OPR) & \num{30.03} \\
  $T_{t4}$ & \SI{1587}{\kelvin} (\num{2857} °R) \\
  Inlet mass flow $W$ & \SI{142.0}{\kilo\gram\per\second} \\
  Fuel--air ratio (FAR) & \num{0.0251} \\
  HPT expansion ratio & \num{3.586} \\
  LPT expansion ratio & \num{4.288} \\
  Total cooling/bleed fraction & \num{0.239} \\
\bottomrule
\end{tabular}

\end{table}

\begin{table}[htbp]
  \centering\small
  \caption{Total conditions at each gas-path station at the design point. Auto-generated.}
  \label{tab:stations}
  % AUTOGENERATED by scripts/make_report_assets.py — do NOT edit by hand
\begin{tabular}{llrr}
\toprule
Station & Location & $P_t$ [bar] & $T_t$ [K] \\
\midrule
  2 & Fan inlet & 0.358 & 245.5 \\
  21 & Fan exit & 0.603 & 289.5 \\
  25 & HPC inlet & 1.149 & 354.2 \\
  3 & HPC exit & 10.739 & 703.9 \\
  4 & Burner exit & 10.159 & 1587.2 \\
  45 & HPT exit & 2.833 & 1141.6 \\
  5 & LPT exit & 0.657 & 806.3 \\
\bottomrule
\end{tabular}

\end{table}

\begin{figure}[htbp]
  \centering
  \includegraphics[width=0.72\textwidth]{design_stations}
  \caption{Total pressure and temperature along the gas path at the design point. Compression
  peaks at station 3 ($\sim$\SI{11}{bar}, \cref{tab:stations}), heat addition at station 4, and
  work extraction across stations 45 and 5.}
  \label{fig:stations}
\end{figure}

The design point converges with a Newton residual below $10^{-6}$ and satisfies all five
acceptance windows defined a priori: BPR $5.10\pm0.2$ (the EEDB target), OPR within its window,
$T_{t4}$ within a plausible band, total cooling-plus-bleed fraction between 18 and 25\,\% of
core flow, and thrust exactly on target (\cref{tab:design}). The resulting TSFC
($\approx0.62$~lb/(lbf$\cdot$h)) lands within 2\,\% of the publicly quoted cruise value for
this engine \emph{without having been used as a calibration target} ---an encouraging
consistency check.

\section{Calibration against measured data (WP1.2)}
\label{sec:calibration}

\subsection{Data sources and philosophy}

Off-design calibration uses exclusively \textbf{measured, certified data} (\cref{tab:tcds}).
The TCDS provides limits and rated thrusts. The ICAO Engine Emissions Databank provides
something less well known: sea-level test-bed \emph{measured fuel flows} at the four
landing-takeoff-cycle thrust settings (100, 85, 30 and 7\,\% of rated thrust $F_{00}$), plus
the engine pressure ratio at rated output (27.61) ---all for the exact -7B26 variant, traceable
to certification testing, and free. Most academic engine models ignore this source.

The anchor points are solved \emph{thrust-matched}: the fuel--air ratio balance is set to hold
$F_n = \mathrm{PC}\cdot F_{00}$ (PC being the power fraction). The model's fuel flow at each
anchor is therefore a genuine \textbf{prediction} confronted with the EEDB measurement ---not a
quantity fitted to it.

\subsection{Calibration iterations}

The first run (model freshly adapted from the example, not yet calibrated to the SLS data)
already predicted takeoff fuel flow within $-2.1$\,\%, but showed two quantitative defects:

\begin{enumerate}
  \item \textbf{N2 above its redline at rated thrust}: \SI{15311}{rpm} against the certified
    limit of \SI{15183}{rpm}. Cause: the reference speed used to scale the generic HPC map
    ---the map's speed-versus-flow shape is not the real compressor's, so the speed labels it
    implies must be anchored deliberately. Fix: HP-shaft design reference from \SI{14324}{} to
    \SI{13940}{rpm}, placing rated-thrust N2 at \SI{14915}{rpm} ($\approx103$\,\%), inside the
    limit.
  \item \textbf{Takeoff pressure ratio 28.99 versus the measured 27.61} ($+5.0$\,\%). Fix: HPC
    design pressure ratio from 9.81 to 9.35. The measured value takes precedence over the
    initial derived estimate, and the cruise OPR becomes a consequence ($\approx30$).
\end{enumerate}

\begin{table}[htbp]
  \centering\small
  \caption{SLS anchors versus the ICAO databank measurements after calibration. $W_f$~model is
  the model's prediction at imposed thrust. Auto-generated.}
  \label{tab:anchors}
  \input{tables/table_anchors}
\end{table}

\begin{figure}[htbp]
  \centering
  \includegraphics[width=0.7\textwidth]{anchor_errors}
  \caption{Fuel-flow prediction error at each anchor against the EEDB measurement. Gray bands
  are the contract tolerances ($\pm3$\,\% at takeoff and climb-out, $\pm5$\,\% at approach and
  idle); the hollow marker (idle) converges but sits outside tolerance and is reported without
  blocking the milestone (\cref{sec:idle}).}
  \label{fig:anchor-errors}
\end{figure}

\begin{figure}[htbp]
  \centering
  \includegraphics[width=0.7\textwidth]{sls_working_line}
  \caption{The model's sea-level-static working line (a continuous sweep of thrust setting)
  against the four measured ICAO databank points. The curvature at low power reflects the
  degraded accuracy of generic maps far from their scaling point.}
  \label{fig:working-line}
\end{figure}

After both corrections (\cref{tab:anchors,fig:anchor-errors,fig:working-line}), the three gated
anchors meet their tolerances: takeoff at $-1.40$\,\% fuel-flow error with the pressure ratio
within $+0.3$\,\% of the measured value and both shaft speeds inside their redlines; climb-out
at $-2.23$\,\%; approach at $-3.08$\,\%.

\subsection{Idle convergence: diagnosis and cure}
\label{sec:idle}

The idle point (7\,\% of $F_{00}$, static) systematically diverges from a cold start. At such
low power the nozzle pressure ratios approach unity ---the exhaust barely pushes--- and the
nozzles' internal static-pressure sub-solver bounces against its lower bound, leaving the
global Newton iteration without a descent direction. Worse, the failed attempt leaves the
point's internal states (flow stations, map positions) corrupted, so re-seeding only the
balance variables does not rescue it: the idle point must \emph{never} start cold at its final
setting. The cure ---anticipated as a generic countermeasure in the specification--- is
\textbf{power continuation}: the point starts at 30\,\% of $F_{00}$ (where it converges from
cold, just like the approach anchor) and walks down warm through
$0.30 \rightarrow 0.20 \rightarrow 0.14 \rightarrow 0.10 \rightarrow 0.08 \rightarrow 0.07$,
re-solving at each step. The logic is encapsulated in the library
(\code{solve\_anchors()}) and covered by the tests.

With continuation, idle converges and predicts fuel flow with $+10.6$\,\% error. This is deep
map-extrapolation territory ---the generic map's speed-flow shape at very low power is not the
real compressor's--- a known limitation of the approach; the anchor is \textbf{reported but not
gated} for milestone H1, a decision recorded in the calibration contract
(\code{gated:~false}).

\section{Baseline decks (WP1.3)}
\label{sec:decks}

Both EHM pipelines need a \emph{baseline}: what a healthy engine reads at any operating
condition, so that measured values can be turned into deviations. The baseline is produced by
sweeping the healthy model over a realistic flight envelope ---altitude blocks with their
plausible Mach ranges (no Mach~0.8 at sea level), ambient temperature offsets of 0, $+15$ and
$+30$\,\si{\degreeCelsius} above ISA, and six thrust settings from 50 to 100\,\% of the locally
available maximum. At each condition the model first finds the maximum thrust at the rating
temperature (point \code{OD\_max}, throttled by $T_{4}$), then solves the part-power point at
the commanded fraction (\code{OD\_pt}) ---the same two-point pattern used by the pyCycle
reference example.

Two engineering lessons from WP1.2 shaped the sweep. First, the walk is ordered for warm
starting within each block (thrust descending), and the independent (altitude, Mach) blocks run
in \emph{parallel processes}, one Problem per worker (project norm N1: long decomposable
computations are parallelized; the sweep takes $3.5$ minutes on all cores versus $10$ minutes
serial). Second ---the lesson learned the hard way--- a failed Newton attempt corrupts the
point's internal states beyond re-seeding, so the recovery action on any failure is to
\emph{rebuild} the problem cold with altitude-appropriate guesses (norm N2); a first version of
the sweep without this recovery lost every point downstream of the first divergence.

The result is stored as a table plus a scattered linear interpolator in the four operating
coordinates, and its quality is measured by leave-one-out cross-validation (each grid point
predicted from all the others), reported per channel in \cref{tab:deck-loo}. Median errors are
a few tenths of a percent; the tail percentiles are dominated by envelope-corner points whose
interpolation simplices span across altitude blocks in raw physical coordinates. This is
expected and temporary: the production baseline of phase F3 interpolates in
\emph{corrected-parameter} space (\cref{sec:stations}), where the $\theta$/$\delta$
normalization collapses most of the altitude and temperature dependence and the same grid is
far denser relative to the remaining variation.

\begin{table}[htbp]
  \centering\small
  \caption{Leave-one-out interpolation error of the baseline deck, per channel.
  Auto-generated.}
  \label{tab:deck-loo}
  % AUTOGENERATED by scripts/make_report_assets.py — do NOT edit by hand
\begin{tabular}{lrrr}
\toprule
Channel & Median [\%] & P95 [\%] & Max [\%] \\
\midrule
  Fuel flow & 0.498 & 15.767 & 90.603 \\
  N1 & 0.268 & 6.369 & 25.721 \\
  N2 & 0.080 & 2.649 & 9.697 \\
  EGT & 0.167 & 6.012 & 24.062 \\
  Net thrust & 0.000 & 16.437 & 85.402 \\
\bottomrule
\end{tabular}

\end{table}

\section{Influence coefficient matrix and observability (WP1.4)}
\label{sec:icm}

\subsection{Health-parameter injection}

To compute the ICM ---and, in phase F2, to simulate degraded engines--- the model must accept
imposed health deviations. pyCycle scales its generic component maps with per-component design
scalars (efficiency and flow); the off-design points normally inherit them from the design
point. In the \code{MPCFM56Study} model those connections are routed through multiplier
components, so that any efficiency or flow-capacity deviation can be dialed in per run. Two
implementation subtleties are worth recording: (i) the multipliers must execute \emph{before}
the off-design point ---the multi-point group runs its children in insertion order, and a
mis-ordered version silently solved the off-design point with unscaled maps; the symptom (a
``healthy'' point that did not reproduce the design point) was caught by the self-consistency
check below; and (ii) holding N1 constant, the CFM56's rating parameter, requires a fuel-flow
balance targeting the LP spool speed through a passthrough component, since the speed is itself
the output of another balance.

The wiring is validated by a self-consistency property that is now a permanent test: the
\emph{healthy} off-design point at cruise with N1 equal to its design value must reproduce the
design point exactly. It does: thrust \SI{5480.0}{lbf}, N2 \SI{13940.0}{rpm}, OPR 30.03, EGT
\SI{1141.6}{\kelvin} ---machine-precision agreement with \cref{tab:design}.

\subsection{The ICM at the reference conditions}

The ICM is computed by central finite differences ($\pm0.5$\,\% per health parameter) at the
two conditions where airline EHM snapshots are taken: cruise (M\,0.78/\SI{35000}{ft}) and
sea-level takeoff, at constant N1. \Cref{fig:icm} shows the cruise matrix for the extended
sensor set; every sign satisfies the physical expectations (e.g.\ a healthier HPT runs cooler
and thriftier: negative EGT and $W_f$ entries in its efficiency column), and these sign checks
run automatically in the test suite.

\begin{figure}[htbp]
  \centering
  \includegraphics[width=\textwidth]{icm_cruise}
  \caption{Influence coefficient matrix at cruise (extended sensor set): \% change of each
  measured channel per $+1$\,\% change of each health parameter. Auto-generated from the
  model.}
  \label{fig:icm}
\end{figure}

\subsection{Observability: the quantitative case for hypothesis H2}

With the cockpit set only (N2, $W_f$, EGT ---N1 is the held rating parameter, so it carries no
information), the ICM has rank 3 for 10 unknowns and the geometry of its columns is damning:
the minimum angle between fault signatures is \textbf{1.3\,degrees}, with seven pairs below
15\,degrees (\cref{tab:confusable}). The most confusable pair, HPC efficiency versus HPT
efficiency, is precisely the classical HPC/HPT ambiguity reported in the GPA
literature~\cite{Li2002} ---here quantified for this engine. Two faults whose signatures differ
by 1.3\,degrees are, at realistic noise levels, indistinguishable to any linear
estimator: this is the smearing mechanism measured, and it defines the ``confusable'' subset on
which hypothesis H2 will compare the AI's isolation ability against WLS-GPA. The extended
sensor set (adding P25/T25/PS3/T3) raises the rank to 7 and multiplies the minimum angle by
roughly five ---the quantitative basis for the sensor-set ablation of contribution C6.

\begin{table}[htbp]
  \centering\small
  \caption{Confusable fault pairs at cruise with cockpit sensors (signature angle
  $<15^{\circ}$). Auto-generated.}
  \label{tab:confusable}
  % AUTOGENERATED by scripts/make_report_assets.py — do NOT edit by hand
\begin{tabular}{llr}
\toprule
Fault A & Fault B & Angle [deg] \\
\midrule
  $\eta_{\mathrm{HPC}}$ & $\eta_{\mathrm{HPT}}$ & 1.32 \\
  $\eta_{\mathrm{FAN}}$ & $\eta_{\mathrm{LPT}}$ & 4.32 \\
  $\eta_{\mathrm{HPT}}$ & $\Gamma_{\mathrm{HPT}}$ & 6.02 \\
  $\Gamma_{\mathrm{FAN}}$ & $\eta_{\mathrm{LPT}}$ & 6.52 \\
  $\eta_{\mathrm{HPC}}$ & $\Gamma_{\mathrm{HPT}}$ & 7.20 \\
  $\eta_{\mathrm{FAN}}$ & $\Gamma_{\mathrm{FAN}}$ & 10.05 \\
  $\eta_{\mathrm{FAN}}$ & $\eta_{\mathrm{LPC}}$ & 13.53 \\
\bottomrule
\end{tabular}

\end{table}

\subsection{Thermodynamics cross-check}

Decision D2 (tabular thermodynamics for speed) is closed by re-solving the design point with
NASA's chemical-equilibrium code (CEA): all key quantities agree within
0.7\,\% (\cref{tab:cea}), an order of magnitude below the $\pm3$\,\% calibration tolerance.
Tabular thermodynamics is therefore retained for all fleet-scale computation.

\begin{table}[htbp]
  \centering\small
  \caption{Design point solved with tabular versus CEA thermodynamics. Auto-generated.}
  \label{tab:cea}
  % AUTOGENERATED by scripts/make_report_assets.py — do NOT edit by hand
\begin{tabular}{lrrr}
\toprule
Quantity & Tabular & CEA & $\Delta$ [\%] \\
\midrule
  Net thrust [lbf] & 5480.0000 & 5480.0000 & -0.000 \\
  TSFC [lb/(lbf$\cdot$h)] & 0.6214 & 0.6249 & +0.553 \\
  Inlet mass flow [lbm/s] & 313.0851 & 315.2029 & +0.676 \\
  Fuel--air ratio & 0.0251 & 0.0250 & -0.122 \\
  HPT expansion ratio & 3.5865 & 3.5902 & +0.104 \\
  LPT expansion ratio & 4.2879 & 4.3109 & +0.537 \\
\bottomrule
\end{tabular}

\end{table}

\section{Declared limitations}
\label{sec:limitations}

The model is \emph{representative}, not certifiable, and its limits are declared for citation
in the results chapters: (i) no transients ---steady-state points only; (ii) variable geometry
(bleed valves, variable stators) implicit in the maps, not modeled; (iii) no Reynolds or
humidity corrections; (iv) absolute N1/N2 labels approximate by construction (generic map
shapes), though respecting the certified redlines at the rated point; (v) the model's EGT is
the HPT-exit temperature (station 45), a consistent proxy ---identical for ground truth and for
both EHM pipelines--- of the certified T49.5, whose displayed value additionally carries the
shunt and trim functions described in \cref{sec:engine}; and (vi) the idle error quoted above.
None of these limitations affects the validity of the AI-versus-traditional comparison, which
rests on the model's sensitivity structure (the ICM), not on its absolute values.

\section{Automated verification}

The chapter's claims are pinned by tests: design-point convergence and acceptance windows;
convergence of all four anchors including the idle continuation; fuel-flow and pressure-ratio
tolerances at the gated anchors; N1/N2 redlines at rated thrust; the Study self-consistency
property (healthy off-design $=$ design); the ICM physical sign checks and its rank-3 cockpit
underdetermination; and the pyCycle environment smoke test. Twelve tests green in continuous
integration at the time of writing.
